\documentclass[12pt]{article}

\usepackage{fouriernc}
\usepackage{amsmath}
\usepackage{enumitem} % alphabetic enumerate

\usepackage{tikz}

\newcommand{\contradiction}{%
\begin{tikzpicture}[rotate=45,x=0.5ex,y=0.5ex]
\draw[line width=.2ex] (0,2) -- (3,2) (0,1) -- (3,1) (1,3) -- (1,0) (2,3) -- (2,0);
\end{tikzpicture}
}

\begin{document}

\section*{Problem-120}
\subsection*{Problem Statement}
Can you factor any other polynomial of the form $a^{2n} + b^{2n}$?

\subsection*{Solution}
From the definition of a polynomial, $n \geq 0$. We shall consider two separate cases: (1) when $n$ is even (2) when $n$ is odd.

\subsubsection*{When $n$ is even, therefore $n$ has the form $2m$ with $m \geq 0$}

\begin{equation*}
	\begin{aligned}
		a^{2n}+b^{2n} &= a^{4m}+b^{4m}\\
					  &= \left( a^{2m} \right)^2 + \left( b^{2m} \right)^2\\
					  &= \left( a^{2m} \right)^2 + 2\cdot a^{2m}\cdot b^{2m}+ \left( b^{2m} \right)^2 - 2\cdot a^{2m}\cdot b^{2m}\\
					  &= \left( a^{2m}+b^{2m} \right)^2 - \left( \sqrt{2} \cdot a^m \cdot b^m \right)^2\\
					  &= \left( a^{2m}+ \sqrt{2} \cdot a^m \cdot b^m + b^{2m} \right)\ \left( a^{2m} - \sqrt{2} \cdot a^m \cdot b^m + b^{2m} \right)
	\end{aligned}
\end{equation*}

\subsubsection*{When $n$ is odd, therefore $n$ has the form $2m+1$ with $m \geq 0$}

Let's consider some examples and try to generalize from them.
\begin{itemize}
\item When $m=0$ therefore $n=1$.
\begin{equation*}
	\begin{aligned}
		a^2+b^2 &= a^2 - \left( -b^2 \right)\\
				&= a^2 - \left( \sqrt{-1}\cdot b \right)^2\\
				&= a^2 - (i\cdot b)^2\\
				&= (a+i\cdot b)\ (a-i\cdot b)
	\end{aligned}
\end{equation*}

\item When $m=1$ therefore $n=3$.
Our polynomial in this case is $a^{6}+b^{6}$. We notice that if we set $a^2 = -b^2$, the polynomial evaluates to zero. So, we might be able to factor the polynomial into $\left(a^2+b^2\right)\ (\ \ldots\ )$. We thus try to extract $a^2+b^2$ from each pair of consecutive terms, introducing temporary terms as necessary.
\begin{equation*}
	\begin{aligned}
		a^6+b^6 &= a^6 + a^4\cdot b^2 - a^4\cdot b^2 - a^2\cdot b^4 + a^2\cdot b^4 + b^6\\
				&= a^4\cdot \left( a^2+b^2 \right) - a^2\cdot b^2\cdot \left( a^2+b^2 \right) + b^4\cdot \left( a^2+b^2 \right)\\
				&= \left( a^2+b^2 \right)\ \left( a^4-a^2\cdot b^2+b^4 \right)
	\end{aligned}
\end{equation*}
\item When $m=2$ therefore $n=5$.
Our polynomial in this case is $a^{10}+b^{10}$. Again, setting $a^2=-b^2$ makes the polynomial zero. So, $a^2+b^2$ is a potential factor. We again try to extract $a^2+b^2$ from each pair of consecutive terms introducing temporary terms as necessary.
\begin{equation*}
	\begin{aligned}
		& a^{10} + b^{10}\\
		&=\ a^{10} + a^8\cdot b^2 - a^8\cdot b^2 - a^6\cdot b^4 + a^6\cdot b^4 + a^4\cdot b^6 - a^4\cdot b^6 - a^2\cdot b^8 + a^2\cdot b^8 + b^{10}\\
		&= a^8 \cdot \left( a^2+b^2 \right) - a^6\cdot b^2\cdot \left( a^2+b^2\right) + a^4\cdot b^4\cdot \left( a^2+b^2 \right) - a^2\cdot b^6\cdot \left( a^2+b^2 \right) + b^8\cdot \left( a^2+b^2 \right)\\
		&= \left( a^2+b^2 \right)\ \left( a^8 - a^6\cdot b^2 + a^4\cdot b^4 - a^2\cdot b^6 + b^8 \right)
	\end{aligned}
\end{equation*}
\end{itemize}

The factoring process for $n > 1$ looks like below:
\begin{equation*}
	\begin{aligned}
	a^{2n}+b^{2n} &= \sum_{k=0}^{n-1}\ (-1)^k\ \left[ a^{2(n-k)} \cdot b^{2k}\ +\ a^{2(n-k-1)}\cdot b^{2(k+1)} \right]\\
				  &= \sum_{k=0}^{n-1}\ (-1)^k\ a^{2(n-k-1)}\cdot b^{2k}\ \left[ a^2+b^2 \right]\\
				  &= \left( a^2+b^2 \right)\ \left( \sum_{k=0}^{n-1}\ (-1)^k\ a^{2(n-k-1)}\cdot b^{2k} \right)
	\end{aligned}
\end{equation*}

From the above three cases, we can generalize the factoring of $a^{2n}+b^{2n}$ when $n$ is odd as follows:
\[
a^{2n}+b^{2n} = \begin{cases} (a+i\cdot b)\ (a-i\cdot b) &\mbox{, when } n=1\\ \left( a^2+b^2 \right)\ \left( \sum_{k=0}^{n-1}\ (-1)^k\ a^{2(n-k-1)}\cdot b^{2k} \right) & \mbox{, otherwise} \end{cases}
\]
\pagebreak
\section*{Problem-122(c)}
\subsection*{Problem Statement}
Factor $a^{10} + a^5 + 1$.
\subsection*{Solution}
\begin{equation*}
	\begin{aligned}
			& a^{10} + a^5 + 1\\
			&= \left( a^5 \right)^2 + 2 \cdot a^5 \cdot \frac{1}{2} + \left( \frac{1}{2} \right)^2 + \frac{3}{4}\\
			&= \left( a^5+\frac{1}{4} \right)^2 + \left( \frac{ \sqrt{3} }{2} \right)^2\\
			&= \left( a^5+\frac{1}{4} \right)^2 - \left( i\cdot \frac{ \sqrt{3} }{2} \right)^2\\
			&= \left( a^5 + \frac{1}{4} + \frac{i\cdot \sqrt{3}}{2} \right)\ \left( a^5 + \frac{1}{4} - \frac{i\cdot \sqrt{3}}{2} \right)\\
			&= \left( a^5 + \frac{1 + 2\sqrt{3}\cdot i}{4} \right)\ \left( a^5 + \frac{1-2\sqrt{3}\cdot i}{4} \right)
	\end{aligned}
\end{equation*}

\section*{Problem-122(d)}
\subsection*{Problem Statement}
Factor $a^3 + b^3 + c^3 -3abc$.
\subsection*{Solution}
We note that setting $a=b=c$ makes the polynomial zero. A possible factor may be $a^2+b^2+c^2-ab-bc-ca$. So, from each set of consecutive six terms we shall extract out this potential factor, introducing temporary terms as necessary.
\begin{equation*}
	\begin{aligned}
		& a^3 + b^3 + c^3 -3abc\\
		&=\ \left(a^3+ab^2+ac^2-a^2b-abc-ca^2\right) + \left( a^2b+b^3+bc^2-ab^2-b^2c-abc \right) \\
		&+ \left( ca^2 + b^2c + c^3 - abc - bc^2 - c^2a \right)\\
		&= a\cdot \left( a^2+b^2+c^2-ab-bc-ca \right) + b\cdot \left( a^2+b^2+c^2-ab-bc-ca \right)\\
		&+\ c\cdot \left( a^2+b^2+c^2-ab-bc-ca \right)\\
		&= (a+b+c)\ \left( a^2+b^2+c^2-ab-bc-ca \right)
	\end{aligned}
\end{equation*}
Yep, it looks like magic, pulling things out of thin air to make it work. But this is expected, since when we multiply factors of a polynomial (the reverse of factoring), we collect the similar terms and in that process some terms may vanish. It reminds me that differentiation is straightforward, integration is hard. Here, multiplication of polynomials is easy, factoring a polynomial is hard.

\section*{Problem-122(e)}
\subsection*{Problem Statement}
Factor $(a+b+c)^3 -a^3-b^3-c^3$.
\subsection*{Solution}
We note that setting $a=-b$ makes the polynomial zero. So, $a+b$ is a potential factor.
\begin{equation*}
	\begin{aligned}
		& (a+b+c)^3 -a^3-b^3-c^3\\
		&=\ \left( 3a^2b+3ab^2 \right) + \left( 3abc+3b^2c \right) + \left( 3a^2c+3abc \right) + \left( 3ac^2 + 3bc^2 \right)\\
		&= 3ab\cdot (a+b) + 3bc\cdot (a+b) + 3ca\cdot (a+b) + 3c^2(a+b)\\
		&= 3(a+b)\ \left(ab+bc+ca+c^2\right)\\
		&= 3(a+b)\ [\ b(a+c)+c(a+c)\ ]\\
		&= 3\ (a+b)\ (b+c)\ (c+a)
	\end{aligned}
\end{equation*}
Actually there is nothing special about the pair $(a,b)$, all three pairs behave identically. We could set $b=-c$ or $c=-a$ to make the polynomial zero. Hence all three $a+b$, $b+c$, and $c+a$ are factors.

\section*{Problem-122(f)}
\subsection*{Problem Statement}
Factor $(a-b)^3 + (b-c)^3 + (c-a)^3$.
\subsection*{Solution}
We note that setting $a=b$ or $b=c$ or $c=a$ makes the polynomial zero. So, possibly the polynomial has the form $k\cdot (a-b)\ (b-c)\ (c-a)$ where $k$ is some constant. Following similar procedure as in $122(e)$ we find below factors:
\[
	(a-b)^3 + (b-c)^3 + (c-a)^3 = 3\ (a-b)\ (b-c)\ (c-a)
\]

\section*{Problem-134}
\subsection*{Problem Statement}
You know that $x + \frac{1}{x}$ is an integer. Prove that $x^n + \frac{1}{x^n}$ is an integer for any $n=1,2,3,$ etc.

\subsection*{Solution}
Let's work with some examples and try to generalize from them.
\begin{itemize}
\item When $n=2$, we need to show $x^2 + \frac{1}{x^2}$ is an integer.
\begin{equation*}
	\begin{aligned}
		\left(x + \frac{1}{x}\right)^2 &= x^2 + 2 \cdot x \cdot \frac{1}{x} + \frac{1}{x^2}\\
		&= x^2+2+\frac{1}{x^2}
	\end{aligned}
\end{equation*}
So, we have $x^2+\frac{1}{x^2} = \left(x + \frac{1}{x}\right)^2-2$. Since $x + \frac{1}{x}$ is an integer, its square is also an integer. And, if we subtract $2$ from the square, the result will still be an integer.

\item When $n=3$, we need to show $x^3 + \frac{1}{x^3}$ is an integer.
\begin{equation*}
	\begin{aligned}
		\left(x + \frac{1}{x}\right)^3 &= x^3 + 3 \cdot x^2 \cdot \frac{1}{x} + 3 \cdot x \cdot \frac{1}{x^2} + \frac{1}{x^3}\\
		&= x^3+\frac{1}{x^3} + 3\cdot \left( x + \frac{1}{x} \right)
	\end{aligned}
\end{equation*}
We have $x^3 + \frac{1}{x^3} = \left(x + \frac{1}{x}\right)^3-3\cdot \left( x + \frac{1}{x} \right)$, which is an integer.

\item When $n=4$, we need to show $x^4 + \frac{1}{x^4}$ is an integer. Using binomial theorem we can expand $\left(x + \frac{1}{x}\right)^4$ as follows:
\begin{equation*}
	\begin{aligned}
		\left(x + \frac{1}{x}\right)^4 &= \sum_{k=0}^{4} \binom{4}{k}\ x^{4-k}\ \left( \frac{1}{x} \right)^k\\
		&= x^4 + \binom{4}{1}x^3 \cdot \frac{1}{x} + \binom{4}{2}x^2\ \cdot \frac{1}{x^2} + \binom{4}{3}x \cdot \frac{1}{x^3} + \frac{1}{x^4}\\
		&= x^4 + \frac{1}{x^4} + \binom{4}{1}\ \left( x^2 + \frac{1}{x^2} \right) + \binom{4}{2}\ \ \ \ \ \ \textrm{// Using}\ \binom{n}{k} = \binom{n}{n-k}
	\end{aligned}
\end{equation*}
We have already shown that $x^2+\frac{1}{x^2}$ is an integer. 

So $x^4 + \frac{1}{x^4} = \left(x + \frac{1}{x}\right)^4 - \binom{4}{1}\ \left( x^2 + \frac{1}{x^2} \right)- \binom{4}{2}$ is an integer.

\item When $n=5$.
\begin{equation*}
	\begin{aligned}
		\left( x + \frac{1}{x} \right)^5 &= \sum_{k=0}^5 \binom{5}{k} x^{5-k} \frac{1}{x^k}\\
		&= x^5 + \frac{1}{x^5} + \binom{5}{1}\left( x^3 + \frac{1}{x^3} \right) + \binom{5}{2} \left( x + \frac{1}{x} \right)
	\end{aligned}
\end{equation*}
\end{itemize}
From the above examples, we see that if we expand $\left(x + \frac{1}{x}\right)^n$ we get terms like $x^m + \frac{1}{x^m}$ where $m \leq n$. This suggests using strong induction. To show $x^n + \frac{1}{x^n}$ is an integer, let's expand $\left(x+\frac{1}{x}\right)^n$.
\begin{equation*}
	\begin{aligned}
		& \left(x+\frac{1}{x}\right)^n\\
		&= \sum_{k=0}^n \binom{n}{k} x^{n-k} \frac{1}{x^{k}}\\
		&= x^n + \frac{1}{x^n} + \sum_{k=1}^{n-1}\binom{n}{k}x^{n-k}\frac{1}{x^k}
	\end{aligned}
\end{equation*}
When $n$ is odd, we can pair up the terms and write as follows:
\[
	\left(x+\frac{1}{x}\right)^n = x^n + \frac{1}{x^n} + \sum_{k=1}^{\frac{n-1}{2}}\binom{n}{k}\left( x^{n-2k} + \frac{1}{x^{n-2k}} \right)
\]
By strong induction we know $x^{n-2k} + \frac{1}{x^{n-2k}}$ are integers, thus $x^n + \frac{1}{x^n}$ is an integer.

When $n$ is even, there are odd number of terms in the binomial expansion, with a lone center at $k = n-k$ which gives an integer $\binom{n}{n/2}$. Thus, for even $n$ we may write:
\[
	\left(x+\frac{1}{x}\right)^n = x^n + \frac{1}{x^n} + \sum_{k=1}^{\frac{n-1}{2}}\binom{n}{k}\left( x^{n-2k} + \frac{1}{x^{n-2k}} \right) + \binom{n}{n/2}
\]

Again, by strong induction we know $x^{n-2k} + \frac{1}{x^{n-2k}}$ are integers, thus $x^n + \frac{1}{x^n}$ is an integer.  $\blacksquare$

\section*{Problem-153}
\subsection*{Problem Statement}
The polynomial $P(x) = x^3 + x^2 -10x + 1$ has three different roots (the authors guarantee it) denoted by $x_1, x_2, x_3$. Write a polynomial with integer coefficients having roots
\begin{enumerate}[label=(\alph*)]
	\item $x_1+1, x_2+1, x_3+1$
	\item $2x_1, 2x_2, 2x_3$
	\item $\frac{1}{x_1}, \frac{1}{x_2}, \frac{1}{x_3}$
\end{enumerate}

\subsection*{Polynomial with roots $x_1+1, x_2+1, x_3+1$}
Since $x_1$ is a root of $P(x)$, $P(x_1) = 0$. Thus, we want to find a polynomial $S(x)$ such that $S(x_1+1) = P(x_1)$ for example.
\begin{equation*}
	\begin{aligned}
		S(x) &= (x-1)^3 + (x-1)^2 - 10\ (x-1) + 1\\
			&= x^3 - 2x^2-9x+11
	\end{aligned}
\end{equation*}

\subsection*{Polynomial with roots $2x_1, 2x_2, 2x_3$}
A good start would be with a polynomial say $T(x)$ such that $T(2x_1) = P(x_1)$.
\begin{equation*}
	\begin{aligned}
		T(x) &= \left( \frac{x}{2} \right)^3 + \left( \frac{x}{2} \right)^2 - 10\ \left( \frac{x}{2} \right) + 1\\
		&= \frac{1}{8}x^3 + \frac{1}{4}x^2-5x+1
	\end{aligned}
\end{equation*}
$T(x)$ is close but not the solution, since we need integer coefficients. So, we take $S(x)$, a constant multiple of $T(x)$, as our solution.
\[
	S(x) = 8\ T(x) = x^3 + 2x^2 -40x + 8
\]

\subsection*{Polynomial with roots $\frac{1}{x_1}, \frac{1}{x_2}, \frac{1}{x_3}$}
Again, a good start would be with a polynomial say $T(x)$ such that 

$T\left( \frac{1}{x_1} \right) = P(x_1)$. Consider $W(x)$ below:

\begin{equation*}
	\begin{aligned}
		W(x) = \frac{1}{x^3} + \frac{1}{x^2} - 10 \cdot \frac{1}{x} + 1
	\end{aligned}
\end{equation*}
Here $W(x)$ is not really a polynomial, since negative power is not allowed in a polynomial by definition. However, $W(x)$ has the property that 

$W\left( \frac{1}{x_1} \right) = P(x_1)$.

To fix that, we can take $S(x) = x^3 \cdot W(x)$ as our solution. Because in that case $S\left( \frac{1}{x_1} \right) = \frac{1}{x_1^3} \cdot W\left( \frac{1}{x_1} \right) = \frac{1}{x_1^3} \cdot P(x_1) = 0$.
\[
	S(x) = x^3 \cdot W(x) = x^3 - 10x^2 + x + 1
\]

\section*{Problem-154}
\subsection*{Problem Statement}
Assume that $x^3 + ax^2 + x + b$ (where $a$ and $b$ are some numbers) is divisible by $x^2-3x+2$. Find $a$ and $b$.

\subsection*{Solution}
Since $P(x) = x^3 + ax^2 + x + b$ is divisible by $x^2-3x+2$, we have:
\begin{equation*}
	\begin{aligned}
		P(x) &= \left( x^2-3x+2 \right) \cdot Q(x)\\
						   &= (x-1) \ (x-2) \ Q(x)
	\end{aligned}
\end{equation*}
So, $1$ and $2$ are roots of $P(x)$.
\begin{equation}
	P(1) = a+b+2 = 0
\end{equation}
\begin{equation}
	P(2) = 4a + b + 10 = 0
\end{equation}
\\
Solving $(1)$ and $(2)$ gives us: $(a, b) = \left( -\frac{8}{3}, \frac{2}{3} \right)$.

\section*{Problem-156}
\subsection*{Problem Statement}
Let us write the values of $P(0),\ P(1),\ P(2),\ \ldots$ for $P(x) = x^2-x-4$.
\[
	-4,\ -4,\ -2,\ 2,\ 8,\ 16,\ 26,\ \ldots
\]
Under any two adjacent numbers write their difference:
\[
\begin{array}{lllllllllllll}
	-4 &   & -4 &   & -2 &   & 2 &   & 8 &   & 16 &    & 26\ \ldots\\
	   & 0 &    & 2 &    & 4 &   & 6 &   & 8 &    & 10\ \ldots &
\end{array}
\]
and repeat the same operation with this sequence of ``first differences'':
\[
\begin{array}{lllllllllllll}
	-4 &   & -4 &   & -2 &   & 2 &   & 8 &   & 16 &    & 26\ \dots\\
	   & 0 &    & 2 &    & 4 &   & 6 &   & 8 &    & 10\ \ldots &   \\
	   &   & \ \ 2 &   &\ \ 2 &   & 2 &   & 2 &   &\ \ 2\ \ldots &    &
\end{array}
\]
Now all numbers are $2$ s. Prove that it is not a coincidence and that all subsequent numbers (called ``second differences'') are also $2$ s.
\subsection*{Solution}
If $P_{\Delta}(x)$ represents ``first differences'', we have:
\begin{equation*}
	\begin{aligned}
		P_{\Delta}(x) &= P(x+1)-P(x)\\
					  &= 2x
	\end{aligned}
\end{equation*}
Let $P_{\Delta \Delta}(x)$ denote the ``second differences''.
\begin{equation*}
	\begin{aligned}
		P_{\Delta \Delta}(x) &= P_{\Delta}(x+1)-P_{\Delta}(x)\\
		                  &= 2
	\end{aligned}
\end{equation*}

\section*{Problem-157}
\subsection*{Problem Statement}
Prove that for any polynomial of degree $2$ all second differences are equal.
\subsection*{Solution}
Let $P(x) = ax^2 + bx + c$ denote a polynomial of degree $2$, where $a, b, c$ are constants. If $P_{\Delta}(x)$ denotes the ``first differences'', we have:
\[
	P_{\Delta}(x) = P(x+1)-P(x) = 2ax + a+b
\]
If $P_{\Delta \Delta}(x)$ represents ``second differences'', we have:
\[
	P_{\Delta \Delta}(x) = P_{\Delta}(x+1) - P_{\Delta}(x) = 2a
\]

\section*{Problem-158}
\subsection*{Problem Statement}
What can be said about polynomials having degree $3$?
\subsection*{Solution}
Third differences are equal.

\section*{Problem-159}
\subsection*{Problem Statement}
(L. Euler) Compute the values $P(x) = x^2+x+41$ for $x=1,\ 2,\ 3,\ \ldots,\ 10$. Check that all these values are prime numbers (having no divisors except $1$ and themselves). Might it be that all of $P(1),\ P(2),\ P(3),\ \ldots$ are prime numbers for this polynomial $P$?
\subsection*{Solution}
Below {\tt Java} code shows that the values of $P(x)$ for $x=1,\ 2,\ 3,\ \ldots,\ 10$ are indeed prime.
\begin{verbatim}
  public static void main(String[] args) {
    for (int x = 1; x <= 10; ++x) {
      if ( isPrime( pOfX(x) ) ) {
        System.out.println( 
          String.format("P(%2d) = %3d is a prime.", x, pOfX(x)) );
      } else {
        System.out.println( 
          String.format("P(%2d) = %3d is NOT a prime.", x, pOfX(x)) );
      }
    }
  }
  private static boolean isPrime(int pval) {
    for (int d = 2; d < pval; ++d) {
      if (pval % d == 0) {
        return false;
      }
    }
    return true;
  }
  private static int pOfX(int x) {
    return x*x + x + 41;
  }
P( 1) =  43 is a prime.
P( 2) =  47 is a prime.
P( 3) =  53 is a prime.
P( 4) =  61 is a prime.
P( 5) =  71 is a prime.
P( 6) =  83 is a prime.
P( 7) =  97 is a prime.
P( 8) = 113 is a prime.
P( 9) = 131 is a prime.
P(10) = 151 is a prime.
\end{verbatim}
However, not all values of $P(x)$ are prime. Consider $P(40)$ for example.

$P(40) = 40^2+40+41 = 41 \cdot 41$.
\\

This is a remarkable polynomial, nonetheless.
\begin{table}[h!]
  \begin{center}
    \begin{tabular}{l|l}
      \textbf{Range of x} & \textbf{P(x) is prime}\\
      \hline
      $1 \leq x \leq 39$ & 100\%\\
      $1 \leq x \leq 100$ & 86\%\\
      $1 \leq x \leq 200$ & 78\%\\
      $1 \leq x \leq 300$ & 70\%\\
      $1 \leq x \leq 400$ & 67.5\%\\
      $1 \leq x \leq 500$ & 65\%\\
      $1 \leq x \leq 600$ & 63.7\%\\
      $1 \leq x \leq 700$ & 61.6\%\\
      $1 \leq x \leq 800$ & 59.8\%\\
      $1 \leq x \leq 900$ & 59\%\\
      $1 \leq x \leq 1000$ & 58.1\%\\
    \end{tabular}
  \end{center}
\end{table}

\section*{Problem-164}
\subsection*{Problem Statement}
Prove that a polynomial of degree not exceeding 2 is defined uniquely by three of its values.

\subsection*{Solution}
\subsubsection*{Defined by three of its values}
A polynomial of degree not exceeding 2 has the general form $P(x) = ax^2 + bx+c$. With three different values for $x_1 \neq x_2 \neq x_3$, we have the below set of three equations:
\setcounter{equation}{0}
\begin{equation}
ax_1^2 + bx_1 + c = \alpha
\end{equation}
\begin{equation}
ax_2^2 + bx_2 + c = \beta
\end{equation}
\begin{equation}
ax_3^2 + bx_3 + c = \gamma
\end{equation}
Solving $(1),\ (2),\ (3)$ gives us the values of $a,\ b,\ c$ and that defines the polynomial.

\subsubsection*{Defined uniquely}
We shall use \textbf{proof by contradiction}. Therefore, we shall assume that there are two distinct polynomials $P(x)$ and $Q(x)$ having degree not exceeding 2 and they agree in three values. This assumption will lead to a contradiction. In other words, for $x_1 \neq x_2 \neq x_3$, we assume $P(x_1) = Q(x_1),\ P(x_2) = Q(x_2),$ and $P(x_3) = Q(x_3)$. The contradiction thus arrived would force us to conclude that $P(x) = Q(x)$.

Consider the polynomial $R(x) = P(x)-Q(x)$. $R(x)$ cannot have degree more than 2 and $R(x_1) = R(x_2) = R(x_3) = 0$. Therefore, $x_1,\ x_2,\ x_3$ are three distinct roots of $R(x)$. So, we can write $R(x) = k\ (x-x_1)\ (x-x_2)\ (x-x_3)$. But then $R(x)$ has degree 3. Contradiction. \contradiction

Therefore, $R(x) = 0$, which leads to $P(x) = Q(x)$.

\end{document}



















